\documentclass[12pt]{article}
\usepackage[utf8]{inputenc}
\usepackage[T1]{fontenc}
\usepackage{hyperref}
\title{Yarawata Language Resources for K-12 Curriculum Development}
\author{Compiled by Claude Sonnet 5.5 (claude-sonnet-5-5) for Andr\'as Kornai}
\date{}
\begin{document}
\maketitle

\section{Introduction}
Yarawata (ISO 639-3 yrw; Glottocode yara1250) is a Trans-New Guinea language of Madang Province, Papua New Guinea. The search results give about 130 speakers (SIL figure from 2003) and a UNESCO rating of Severely Endangered. Glottolog places it in the Numugenan branch of the Northern Adelbert group. Joshua Project gives a community population of about 300 and describes the group as relatively isolated. My searches found no school materials in Yarawata.

\section{Primary Educational Resources}
No primers, textbooks or school programmes for Yarawata turned up in search. Nothing here can be used directly in a classroom. A curriculum would have to be built from scratch with the community.

\section{Linguistic and Reference Materials}
\begin{itemize}
\item Wikipedia overview of the language: \url{https://en.wikipedia.org/wiki/Yarawata_language}
\item ASJP wordlist, a short comparative list and not a dictionary: \url{https://asjp.clld.org/languages/YARAWATA}
\item Glottolog classification (this is the Glottolog map page for yara1250): \url{https://glottolog.org/resource/languoid/id/yara1250.bigmap.html}
\end{itemize}

\section{Religious and Community Resources}
\begin{itemize}
\item Joshua Project language page: \url{https://joshuaproject.net/languages/yrw}
\item Joshua Project people-group page, which says most Yarawata consider themselves Christian: \url{https://joshuaproject.net/index.php/people_groups/15993}
\item Scripture Earth lists the language, but Joshua Project reports no completed Bible or New Testament, so the actual scripture holdings are unclear: \url{https://www.scriptureearth.org/00eng.php?iso=yrw}
\end{itemize}

\section{Digital and Audio Resources}
\begin{itemize}
\item Endangered Languages Project entry: \url{https://endangeredlanguages.com/lang/3165}
\end{itemize}
I found no audio or video collections.

\section{Curriculum Development Context}
Very little documentation is available online: a short wordlist and catalogue entries. A K-12 curriculum would depend on fieldwork with the roughly 130 speakers, who are probably mostly older adults. Work should be led by the community and begin with documentation and orthography. The status of any Bible translation should be checked directly with Scripture Earth or SIL.
\end{document}